%=========================================================================
%  Tong Chen — One-page résumé (condensed from tong_chen_cv.tex)
%  Compiles with pdfLaTeX. Keep to a single page.
%=========================================================================
\documentclass[11pt]{article}

\usepackage[margin=0.75in]{geometry}
\usepackage{enumitem}
\usepackage{titlesec}
\usepackage[hidelinks]{hyperref}
\usepackage{xcolor}

% --- Section styling: small-caps heading with a rule underneath ---
\titleformat{\section}{\large\bfseries\scshape}{}{0em}{}[\vspace{-4pt}\titlerule]
\titlespacing{\section}{0pt}{8pt}{4pt}

% --- List defaults ---
\setlist[itemize]{leftmargin=1.4em, itemsep=1pt, topsep=1pt, parsep=0pt}
\setlist[enumerate]{leftmargin=1.6em, itemsep=1pt, topsep=1pt, parsep=0pt}

\setlength{\parindent}{0pt}
\pagestyle{empty}

% --- Helper: bold left, right-aligned date ---
\newcommand{\dated}[2]{\textbf{#1}\hfill #2\par}

\begin{document}

%========================= HEADER =========================
\begin{center}
  {\LARGE\textbf{Tong Chen}}\\[4pt]
  470-641-4428 \,$\cdot$\, \href{mailto:txc5603@mavs.uta.edu}{txc5603@mavs.uta.edu}\\[2pt]
  \href{https://tongchen2010.github.io}{tongchen2010.github.io} \,$\cdot$\,
  \href{https://github.com/tongchen2010}{github.com/tongchen2010} \,$\cdot$\,
  \href{https://www.linkedin.com/in/tong-chen-850935124}{linkedin.com/in/tong-chen-850935124}
\end{center}

%========================= RESEARCH INTERESTS =========================
\section{Research Interests}
Machine learning and AI for medical imaging, with a focus on Alzheimer's disease (AD) and
Lewy body dementia (LBD): self-supervised representation learning for longitudinal brain MRI
and disease-progression modeling across the cognitive impairment continuum.

%========================= EDUCATION =========================
\section{Education}
\dated{Ph.D. in Computer Science, The University of Texas at Arlington}{Aug 2023 -- Present}
\textit{Advisor: Dr. Dajiang Zhu \quad (expected Dec 2026)}\par
\vspace{2pt}
\dated{M.S. in Computer Science, Kennesaw State University}{Jan 2021 -- Dec 2022}
\vspace{2pt}
\dated{B.S. in Computer Science, Kennesaw State University}{Aug 2012 -- May 2018}

%========================= SELECTED PUBLICATIONS =========================
\section{Selected Publications}
{\small (\textbf{Bold} = author; * = equal contribution. Full list at
\href{https://tongchen2010.github.io/publications/}{tongchen2010.github.io}.)}
\begin{enumerate}
  \item Minheng Chen, \textbf{Tong Chen}, Yan Zhuang, et al. ``Community-Level Modeling of
        Gyral Folding Patterns for Robust and Anatomically Informed Individualized Brain
        Mapping.'' \textit{NeuroImage}, 2026.
  \item \textbf{Tong Chen}, Minheng Chen, Jing Zhang, et al. ``HiLoGraph: Hierarchical-
        Longitudinal Brain Network Representation Learning.'' \textit{MICCAI}, 2026.
        \textbf{(Early accept, top 9\%)}
  \item \textbf{Tong Chen}, Minheng Chen, Jing Zhang, et al. ``A Unified Continuous Staging
        Framework for Alzheimer's Disease and Lewy Body Dementia via Hierarchical Anatomical
        Features.'' \textit{MICCAI}, 2025. \textbf{(Early accept, top 9\%)}
  \item Chao Cao*, \textbf{Tong Chen}*, Nan Zhao*, et al. ``Gyral-Sulcal-Net: An Integrated
        Network Representation of Brain Folding Patterns.'' \textit{IEEE ISBI}, 2026.
        \textbf{(Co-first author)}
\end{enumerate}

%========================= RESEARCH EXPERIENCE =========================
\section{Research Experience}
\dated{Graduate Research Assistant, UT Arlington (Advisor: Dr. Dajiang Zhu)}{Aug 2023 -- Present}
\begin{itemize}
  \item Developed graph neural network methods for multi-class dementia classification and
        continuous disease-progression modeling on large multi-cohort neuroimaging data.
  \item Built sMRI and DTI preprocessing and analysis pipelines for multi-cohort studies;
        first-author work at MICCAI and AAIC and a co-authored \textit{NeuroImage} article.
\end{itemize}
\vspace{2pt}
\dated{Research Assistant, UNC Chapel Hill (Advisor: Dr. Gang Li)}{May 2025 -- Aug 2025}
\begin{itemize}
  \item Developed cortical surface-based machine learning methods on Human Connectome
        Project (HCP) data.
\end{itemize}

%========================= TECHNICAL SKILLS =========================
\section{Technical Skills}
\textbf{Programming:} Python, MATLAB \quad \textbf{Frameworks:} PyTorch, PyTorch Geometric, scikit-learn \\
\textbf{Models:} Graph Neural Networks, Transformers, State-Space Models (Mamba), Multimodal LLMs \\
\textbf{Neuroimaging:} FreeSurfer, FSL, SPM12, CAT12, Connectome Workbench, DSI Studio, MRtrix3 \\
\textbf{Datasets:} ADNI, NACC, AABC, HCP, BCP, in-house LBD cohorts

\end{document}
